%% !TEX root = ../manual_gaolv5.tex
% Chapter 13.
%% Chapter 13: Floating-point numbers, errors and debugging. Sources:
%% gaol.tex of f5bf180 lines 6007-6123 and 6156-6404 (the rounding functions
%% of the FPU, lines 6124-6155, are in Chapter 4), doc/using.md lines 931-998.
\chapter{Floating-point numbers, errors and debugging}%
\label{chap:errors}\label{chap:numbers}

An interval library rests on the floating-point numbers of its machine. The
bounds of the intervals of GAOL are \code{double} values, and a program that
builds an interval bound by bound, or that examines the result of a
computation, needs a few operations on these values: the neighbours of a
double, the number of doubles between two of them, a minimum and a maximum
that tell the two zeros apart. The first part of this chapter describes the
constants and the functions on \code{double} that GAOL defines beside its
intervals, the constants that tell a program the version of the headers it
was compiled with, and two utility functions on integers.

\medskip
The second part is about the errors. GAOL reports an error in one of two
ways, by throwing an exception or by printing a message and aborting, and the
way depends on how the library was built. GAOL v5 throws the exceptions of
the standard library, where GAOL~4 had classes of its own, so that a handler
that knows nothing of GAOL says what went wrong.\newinvfive{} We describe
the mechanism, the macro with which a program adds its own errors, the two
exceptions that GAOL throws and the text of their messages, and the warnings.
The last part presents the debugging facilities of the Debug builds.

\medskip
The chapter is organized as follows.
Section~\ref{sec:floating-point-constants} gives the floating-point constants
and Section~\ref{sec:fp-functions} the floating-point functions.
Section~\ref{sec:version} describes the version information and
Section~\ref{sec:additional-functions} the additional functions.
Section~\ref{sec:errors-handling} presents the error handling, the exceptions
and the warnings, and Section~\ref{sec:debugging-facilities} the debugging
facilities. Section~\ref{sec:errors-discussion} discusses the choices made.

\section{Floating-point constants}\label{sec:floating-point-constants}
\cindex{floating-point constants}\cindex{NaN}\cindex{infinity}%
\findex{GAOL\_NAN}\findex{GAOL\_INFINITY}%

In addition to the constants available through
\code{numeric\_limits<double>}, GAOL defines the two macros of the table below, which expand into \code{double} values.

\begin{center}
  \renewcommand{\arraystretch}{1.1}
  \begin{tabular}{lc}
    \hline
    \multicolumn{1}{c}{\bfseries Constant (\texttt{double})} &
    \multicolumn{1}{c}{\bfseries Value}\\
    \hline\hline
    \element{GAOL\_NAN} & NaN (quiet)\\
    \element{GAOL\_INFINITY} & $+\infty$ \\
    \hline
  \end{tabular}
\end{center}

\begin{Rem}
GAOL~4 also defined the doubles \code{two\_pi}, \code{pi}, \code{half\_pi},
\code{pi\_dn}, \code{pi\_up}, \code{half\_pi\_dn}, \code{half\_pi\_up},
\code{ln2\_dn}, \code{ln2\_up}, \code{two\_power\_51} and
\code{two\_power\_52}, in the namespace of GAOL. A program that opened it with
\code{using namespace gaol} found them among its own names, and a \code{pi}
of its own was ambiguous. GAOL v5 no longer declares them for the code that
uses it, apart from the four doubles around $\pi$ and $\pi/2$, which
\code{gaol} and \code{gaol\_ieee1788} name \code{gaol\_pi\_dn},
\code{gaol\_pi\_up}, \code{gaol\_half\_pi\_dn} and
\code{gaol\_half\_pi\_up}.\newinvfive{} The bounds of $\pi$, $\pi/2$ and $2\pi$ are those of
\code{interval::pi()}, \code{interval::half\_pi()} and
\code{interval::two\_pi()} (\pxref{sec:interval-constants}),
\code{interval::pi().left()} being \roundDn{\pi}; the bounds of $\ln 2$ are
those of \code{log(interval(2.0))}.
\end{Rem}

\section{Floating-point functions}\label{sec:fp-functions}
\cindex{floating-point functions}%

GAOL provides the following functions on \code{double} values. The functions
that set the rounding direction of the floating-point unit are described with
the rounding direction (\pxref{chap:fpu}).

\begin{deffun}{bool}{feven}{(double \var{x})}
  Returns \code{true} whenever \var{x} is even. The function is not meant for
  an infinity or a NaN.

  \begin{example}
cout << boolalpha << feven(3.0) << ' ' << feven(4.0) << ' ' << feven(3.5);
@outputs^false true false~
  \end{example}
\end{deffun}

\begin{deffun}{double}{next\_float}{(double \var{x})}
  Returns the smallest \code{double} greater than \var{x}, the neighbour
  \nextdouble{x} of \var{x}; \code{next\_float(+oo)} is $+\infty$, and \code{next\_float(DBL\_MAX)} is $+\infty$ with every C library.
\end{deffun}

\begin{deffun}{double}{previous\_float}{(double \var{x})}
  Returns the greatest \code{double} smaller than \var{x}, the neighbour
  \prevdouble{x} of \var{x}.
\end{deffun}

\begin{deffun}{bool}{is\_signed}{(double \var{x})}
  Returns \code{true} whenever the sign bit of \var{x} is set: for the negative numbers and $-0$. With \code{using namespace std;} and \code{using namespace gaol;}, as the examples of this manual assume, the name is ambiguous with the class template \code{std::is\_signed}: the call is written \code{gaol::is\_signed(x)}. No provision is made for a
  NaN: a program that wants to test for the negative numbers and $-0$ alone
  has to test also whether \var{x} is a NaN, with \code{std::isnan()}.
\end{deffun}

\begin{deffun}{double}{minimum}{(double \var{x}, double \var{y})}
  Returns the minimum \code{double} value of \var{x} and \var{y}. The
  function is commutative and returns $-0$ when comparing $-0$ and $+0$:
\begin{equation*}
\left\{\begin{array}{lll}
  \min(x,y) &=\min(y,x), &\forall x\neq\text{NaN}, \forall y\neq\text{NaN}\\
  \min(x,\text{NaN}) &=\min(\text{NaN},x)=\text{NaN}, &\forall x\\
  \min(-0,0) &=\min(0,-0)=-0
\end{array}\right.
\end{equation*}
\end{deffun}

\begin{deffun}{double}{maximum}{(double \var{x},double \var{y})}
  Returns the maximum \code{double} value of \var{x} and \var{y}. The
  function is commutative and returns $+0$ when comparing $-0$ and $+0$:
\begin{equation*}
\left\{\begin{array}{lll}
  \max(x,y) &= \max(y,x), &\forall x\neq\text{NaN}, \forall y\neq\text{NaN} \\
  \max(x,\text{NaN})&=\max(\text{NaN},x)=\text{NaN}, &\forall x\\
  \max(-0,0)&=\max(0,-0)=0
\end{array}\right.
\end{equation*}
\end{deffun}

\begin{deffun}{unsigned long long}{nb\_fp\_numbers}{(double \var{a}, double \var{b})}
  Returns the number of floating-point numbers in the interval
  \itv{\var{a}}{\var{b}}. In particular:
  \begin{itemize}
  \item \code{nb\_fp\_numbers(a,next\_float(a)) == 2};
  \item \code{nb\_fp\_numbers(a,a) == 1}.
  \end{itemize}
  The two zeros are the same number: \code{nb\_fp\_numbers(-0.0,0.0)} is 1,
  and \code{nb\_fp\_numbers(-0.0,1.0)}, \code{nb\_fp\_numbers(-1.0,0.0)} and
  \code{nb\_fp\_numbers(0.0,1.0)} are all $4607182418800017409$ (GAOL~4
  returned $13830554455654793217$ for the first two).\newinvfive

  The precondition is that \var{a} be lower than or equal to \var{b}.

  The result is an \code{unsigned long long}. GAOL~4 named its type with the
  macro \code{ULONGLONGINT}, which GAOL v5 no longer defines, as it prefixes
  with \code{GAOL\_} the macros of its headers.\newinvfive

  The function throws a \code{std::domain\_error} exception
  (\pxref{sec:exceptions}) if either \var{a} or \var{b} is a NaN or an
  infinity, or if \var{a} is greater than \var{b}; where the exceptions are
  disabled, it calls \code{gaol\_error} and aborts. It returns no number
  then, where the earlier editions of this manual gave
  \code{numeric\_limits<unsigned long long>::max()}.\newinvfive
\begin{example}
cout << nb_fp_numbers(1.0, next_float(1.0)) << ' '
     << nb_fp_numbers(1.0, 2.0);
@outputs^2 4503599627370497~
\end{example}
\end{deffun}

\section{Version information}\label{sec:version}
\cindex{version of the library}%

The headers of the library define four constants that give the version of
GAOL a program was compiled with, from the macros \code{GAOL\_MAJOR\_VERSION}, \code{GAOL\_MINOR\_VERSION}, \code{GAOL\_MICRO\_VERSION} and \code{GAOL\_VERSION} of \file{gaol/gaol\_configuration.h}, which a program tests with \code{\#if} (\pxref{sec:install-configuration-header}). The
versioning scheme is the one of the Apache Software Foundation, described at
\url{http://apr.apache.org/versioning.html}. The constants are declared in the namespace \code{gaol\_core} and named in
\code{gaol} and \code{gaol\_ieee1788}.\newinvfive

\begin{defconst}{unsigned int}{version\_major}
  Major version of the library: 5 for GAOL v5.
\end{defconst}

\begin{defconst}{unsigned int}{version\_minor}
  Minor version of the library.
\end{defconst}

\begin{defconst}{unsigned int}{version\_micro}
  Patch version of the library.
\end{defconst}

\begin{defconst}{const char *const}{version}
  Version of the library as a string.
\begin{example}
cout << gaol::version;
@outputs^@version~
\end{example}
\end{defconst}

\section{Additional functions}\label{sec:additional-functions}

The following functions are utility functions, not necessarily related to
intervals or to floating-point numbers.

\begin{defTfun}{template <typename T>}{bool}{odd}{(const \var{T}\& \var{x})}
  Returns \code{true} if \var{x} is odd and \code{false} otherwise. The type
  \var{T} may be any type that provides the operator \code{\&} (``bitwise
  \code{and}'') with the same semantics as the one of the \code{int}s.
\end{defTfun}

\begin{defTfun}{template <typename T>}{bool}{even}{(const \var{T}\& \var{x})}
  Returns \code{true} if \var{x} is even and \code{false} otherwise. The type
  \var{T} may be any type that provides the operator \code{\&} (``bitwise
  \code{and}'') with the same semantics as the one of the \code{int}s.
\end{defTfun}

\section{Error handling}\label{sec:errors-handling}
\cindex{errors}\cindex{error handling}%

A program that uses GAOL may see an error reported in two different ways:
\begin{itemize}
\item[$\bullet$] \textbf{By an exception.} GAOL throws an exception of the
  standard library, which the program may catch.
\item[$\bullet$] \textbf{By a message.} GAOL prints a message on the
  standard error, and the program aborts.
\end{itemize}

\noindent
The exceptions are those of the standard library, declared in
\code{<stdexcept>}. GAOL~4 had classes of its own, \code{gaol\_exception},
\code{input\_format\_error}, \code{invalid\_action\_error} and
\code{unavailable\_feature\_error}, which GAOL v5 no longer has
(\pxref{sec:exceptions}).\newinvfive

\medskip
The mechanism in use depends on the way the library is configured
(Chapter~\ref{chap:installation}). With the option
\loption{enable-exceptions=yes} of \cmd{configure}, the default, every error of an operation
is reported by an exception; an AVX-512 build started on a processor without these instructions prints a message and aborts in every build. A build where
\code{GAOL\_EXCEPTIONS\_ENABLED} is 0, obtained with \loption{disable-exceptions}
of \cmd{configure}, \option{-Denable-exception=false} of meson or
\option{-DGAOL\_EXCEPTIONS\_ENABLED=OFF} of CMake, prints a
message on the standard error and aborts instead of throwing. This is a setting
of the developers, for the tests of that path. Relying on exceptions is more in the
spirit of C++, though it may incur some overhead.

\begin{Rem}
In a build that aborts, the reader of strings is an exception to this rule. An error it
finds as it computes, a wrong call (\code{nth\_root(8, 1.5)},
\code{sin(1, 2)}), an uncertain number as a bound, two different bounds in
\code{<a, b>} or a bound of it that is no double make it fail as a syntax
error does: \code{parse\_interval()} returns \code{false}, and
\code{gaol\_ieee1788::textToInterval()} the empty set, as IEEE 1788-2015 has
it (12.12.7).
\code{gaol::textToInterval()} and the operator \code{>{}>} print a message
and abort then, as for any string that is no interval.\newinvfive
\end{Rem}

\medskip
Where memory runs out, GAOL throws the \code{std::bad\_alloc} of C++. The
reader of strings does too, and frees what it allocated; in GAOL~4, flex ended the program with \code{exit(2)} when it could not allocate its buffer, and the parser read a stack it could not grow as a string nested too deep.\newinvfive

\medskip
It is up to the user to comply with this mechanism when adding error
reporting code to a program. GAOL defines the following macro, to be used
whenever one wants to report an error.

\begin{defmacroA}{gaol\_ERROR}{(\var{excep},\var{msg})}
  The behaviour of the macro depends on the value chosen for the option
  \loption{enable-exceptions}. If the exceptions are enabled, an exception of
  the class \var{excep} is thrown with the text
  ``\code{file:line: }\var{msg}'', \code{file} and \code{line} being those of
  the place where the macro is written (\code{\_\_FILE\_\_}, cut after its last
  \code{/} or \code{\textbackslash}, and \code{\_\_LINE\_\_}; \code{error\_text()}
  of \file{gaol/gaol\_common.h}): the path of the machine that built GAOL is not
  in the message. The class \var{excep} derives from \code{std::exception},
  in practice one of \code{<stdexcept>}; the compiler refuses any other type,
  by a \code{static\_assert}, in a build that aborts as in one that throws.
  Otherwise, the program aborts with the message \var{msg}.\newinvfive

\begin{example}
interval x;

@emph^[Code manipulating x]~

if (x.is_empty()) {
   gaol_ERROR(std::domain_error,"Emptiness of one interval");
}
\end{example}
\end{defmacroA}

The function \code{gaol\_error()} is defined as follows.

\deffunx{void}{gaol\_error}{(const char *\var{err})}
\begin{deffun}{void}{gaol\_error}{(const char *\var{file}, int \var{line},
const char *\var{err})}
Displays a message on the standard error output. The ternary version should
be called with the macro \code{GAOL\_FILE\_POS} for its first two parameters
(\pxref{sec:exceptions}). A null \var{file} is written
\code{(unknown file)}, and a null \var{err} is an empty text.\newinvfive
\end{deffun}

\subsection{Exceptions}\label{sec:exceptions}
\cindex{exceptions}\findex{std::invalid\_argument}\findex{std::domain\_error}%
\findex{std::logic\_error}\findex{std::bad\_alloc}%

GAOL throws the exceptions of the standard library, and none of its
own.\newinvfive{} It throws two of them to report an error:
\begin{itemize}
\item \code{std::invalid\_argument}, when a string is no interval
  (\pxref{sec:invalid_argument});
\item \code{std::domain\_error}, when a function is called with an argument
  it does not take (\pxref{sec:domain_error}), as \code{bisect()} is for a
  ratio outside $(0, 1)$.
\end{itemize}

\noindent
Both derive from \code{std::logic\_error}, which derives from
\code{std::exception}. A single handler of \code{std::logic\_error} catches
both, but the handler of one leaves the other alone, and none of them
catches the \code{std::bad\_alloc} that GAOL throws where memory runs out.

\medskip
The \code{what()} of these exceptions is
``\code{file:line: }\var{explanation}'': the name of the source file of GAOL,
without its directories, and the line that threw the exception, then what
went wrong, as in \code{gaol\_interval.cpp:48: Syntax error in interval
initialization: [1, 2}. For the errors that the reader of strings finds in a
call (\code{nth\_root(8, 1.5)}, \code{sin(1, 2)}), the file is that of the
grammar, \code{gaol\_interval\_parser.ypp}, and the line one of that file,
which the directives \code{\#line} of the generated parser give. A handler
of \code{std::exception}, and the message of an exception that nothing
catches, say what went wrong, where GAOL~4 gave \code{std::exception}
whatever the error. The text is valid as long as the exception is.

\begin{Ex}
The program below gives the reader of strings a text that is no interval,
and catches the exception as a \code{std::exception}, knowing nothing of
GAOL.
\begin{example}
try {
  interval x = gaol::textToInterval("[1, 2");
} catch (const std::exception& e) {
  std::cerr << e.what() << '\n';
}
\end{example}
The handler writes \code{gaol\_interval.cpp:48: Syntax error in interval
initialization: [1, 2} on the standard error. Had nothing caught the
exception, the program would have ended with the same text; with libstdc++,
the output is the following.
\begin{onscreen}
terminate called after throwing an instance of 'std::invalid_argument'
  what():  gaol_interval.cpp:48: Syntax error in interval initialization: [1, 2
\end{onscreen}
\end{Ex}

\begin{defmacro}{GAOL\_FILE\_POS}
Expands into the first two arguments of \code{gaol\_error()} and
\code{gaol\_warning()}, the name of the file and the line:

\begin{example}
if (@emph^[some condition]~) {
  gaol_error(GAOL_FILE_POS, "No additional information");
}
\end{example}
\end{defmacro}

\begin{Rem}
What GAOL~4 had is gone: the base class \code{gaol\_exception}, with its
accessors \code{file()}, \code{line()} and \code{explanation()} and the
operator ``\code{<{}<}'' that wrote them, and the classes
\code{input\_format\_error}, \code{invalid\_action\_error} and
\code{unavailable\_feature\_error} of \file{gaol/gaol\_exceptions.h}. The
classes become those of the standard library:
\begin{itemize}
\item \code{gaol::input\_format\_error} becomes
  \code{std::invalid\_argument};
\item \code{gaol::invalid\_action\_error} becomes
  \code{std::domain\_error}.
\end{itemize}
\noindent
What \code{explanation()} gave follows the file and the line in
\code{what()}. \code{unavailable\_feature\_error} has no replacement: \code{atan2()}, which threw it in GAOL~4, is implemented.\newinvfive
\end{Rem}

\subsection{The \code{std::invalid\_argument} exception}%
\label{sec:invalid_argument}

The exception \code{std::invalid\_argument} is thrown whenever a program
attempts to create an interval from an invalid string. This situation may
occur when reading an interval from a stream with the operator \code{>{}>},
or when reading a string with \code{gaol::textToInterval()}, from one string
or from two, a null pointer being no string (\pxref{sec:reading-intervals}).
The example prints what follows the first colon and the blank of
\code{what()}, the explanation, which the file and the line precede.

\begin{example}
try {
  interval x = textToInterval("sin(1)+");
} catch (const std::invalid_argument& e) {
  string text = e.what();
  cout << text.substr(text.find(": ") + 2);
}
@outputs^Syntax error in interval initialization: sin(1)+~
\end{example}

\subsection{The \code{std::domain\_error} exception}\label{sec:domain_error}

The exception \code{std::domain\_error} is thrown in three cases only:
\code{nb\_fp\_numbers()} with NaNs as parameters (\pxref{sec:fp-functions}),
\code{bisect()} with a ratio outside $(0, 1)$ (\pxref{sec:splitting}), and a
string that gives a function an argument it does not take, as
\code{"nth\_root(8, 1.5)"} does.

\subsection{Warnings}\label{sec:warnings}
\cindex{warnings}%

GAOL writes no warning itself; the following function writes one, as
\code{gaol\_error()} writes an error, on the standard error output
\code{std::cerr}, for the code that uses GAOL.

\deffunx{void}{gaol\_warning}{(const char *\var{warn})}
\begin{deffun}{void}{gaol\_warning}{(const char *\var{file}, int \var{line}, const char *\var{warn})}
Prints the message \var{warn} on the standard error output. The second form
should be called with the macro \code{GAOL\_FILE\_POS} for its first two
parameters. A null \var{file} is written \code{(unknown file)}, and a null
\var{warn} is an empty text.\newinvfive
\end{deffun}

\section{Debugging facilities}\label{sec:debugging-facilities}
\cindex{debugging}\findex{GAOL\_DEBUGGING}\findex{NDEBUG}%

The debugging facilities described hereunder are available only when GAOL
has been configured with them. The Debug build of each build enables them:
\option{-DCMAKE\_BUILD\_TYPE=Debug} with CMake, the option
\loption{enable-debug} of \cmd{configure} (\pxref{sec:configuration}), or
\loption{buildtype=debug} with meson; \option{-DGAOL\_DEBUG=ON} with CMake
and \option{-Denable-debug=true} with meson enable them whatever the build
type.\newinvfive{} These builds define \code{GAOL\_DEBUGGING} to 1 in
\file{gaol/gaol\_configuration.h} (with a generator of several configurations, Visual Studio or Xcode, only for the sources of GAOL in the configuration Debug: the header keeps 0, and \code{GAOL\_DEBUG\_VERBOSE} expands to nothing in the program), where it is 0 otherwise, as every option
of that file, so that it is read with \code{\#if} (\code{\#ifdef} is true of
it in every build). The Debug builds compile without \code{NDEBUG} too, so
that \code{assert()} stays active.\newinvfive

\begin{deffun}{int}{debug\_level}{}
  Global variable that remembers the current value of the debugging level.
  The variable is set when the library is initialized, by \code{init()}
  (\pxref{chap:init}). It is declared in the namespace \code{gaol\_core},
  and named in \code{gaol} too.
\end{deffun}

\begin{defmacroA}{GAOL\_DEBUG\_VERBOSE}{(\var{lvl},\var{cmd})}
  Executes \var{cmd} if \var{lvl} is lower than or equal to the current
  debugging level, the variable \code{debug\_level} above, which the macro names unqualified: it is used where \code{gaol} or \code{gaol\_core} is open. The macro was
  named \code{GAOL\_DEBUG} before GAOL v5; this name is now the option of
  CMake that builds GAOL for debugging.\newinvfive

  The macro expands to nothing when the library was not configured with the
  debugging facilities.
\end{defmacroA}

\noindent
A possible use of this macro is the following:

\begin{example}
interval x(-10,10);

@emph^[Some code]~

GAOL_DEBUG_VERBOSE(1,cout << "The value of x is " << x);
x += interval(3.5,4.5);
GAOL_DEBUG_VERBOSE(2,cout << "Now the value of x is " << x);
\end{example}

\noindent
The first message is displayed whenever the debugging level is at least 1,
and the second one only if it is at least 2.

\section{Discussion}\label{sec:errors-discussion}

This chapter gathered what GAOL v5 offers around its intervals. The
constants and the functions on \code{double} let a program reason on the
bounds themselves: the neighbours of a double, the number of doubles an
interval holds, where the two zeros count as one, and a minimum and a
maximum that order $-0$ before $+0$ and propagate a NaN. The doubles of
GAOL~4 that named $\pi$ and $\ln 2$ left the namespace, so that a program
opening it keeps a \code{pi} of its own; the constants of the class
\code{interval} replace them.\newinvfive{} Four constants give the version
of the headers a program was compiled with.

\medskip
On the errors, GAOL v5 made one choice: it throws the exceptions of the
standard library, \code{std::invalid\_argument} for a string that is no
interval and \code{std::domain\_error} for an argument a function does not
take, and none of its own.\newinvfive{} A handler that knows nothing of GAOL
reads in \code{what()} the file and the line of GAOL and what went wrong,
and a program written for GAOL~4 changes the classes it catches
(Chapter~\ref{chap:changes}). The build that aborts instead of throwing
remains available in the three builds, as a setting of the developers
(\option{-DGAOL\_EXCEPTIONS\_ENABLED=OFF} with CMake); in that build, an error that the reader of strings finds
makes the reading fail, so that \code{gaol\_ieee1788::textToInterval()}
returns the empty set as IEEE 1788-2015 has it, rather than end the program.
The debugging facilities belong to the Debug builds, and
\code{GAOL\_DEBUG\_VERBOSE} expands to nothing elsewhere.

\medskip
The next chapter describes the tests of GAOL v5 and its continuous
integration, which exercise the exceptions of this chapter among the rest of
the library, on every platform GAOL supports.
